%!TeX TxS-program:compile =txs:///pythontex
\documentclass{mwart}
\usepackage[utf8]{inputenc}
\usepackage{polski}
\usepackage{pythontex}
\usepackage{graphicx}
\usepackage{pgf}
\usepackage{lmodern,microtype}
\usepackage{amsfonts}
\usepackage{float}

\addtolength{\topmargin}{-1cm}
\addtolength{\oddsidemargin}{-2cm}
\addtolength{\textwidth}{4cm}
\addtolength{\textheight}{0cm}

\def \R{\mathbb{R}}

\author{Marta Hałas 282264}
\title{\textbf{Lista 4}}

\begin{document}
	\maketitle
	\medskip
	
	
	\section*{\textbf{Zadanie 2}}\label{Zadanie 2}

	
	\begin{figure}[H]
		\begin{center}
			\includegraphics[width=15cm]{esaaaaaaaaaaaa.png}
			\caption{Wykres: Porównanie czasów działania algorytmów sortowania}
		\end{center}
	\end{figure}

	\medskip

	\begin{pysub}
	 Najwolniejszy okazał się algorytm Sortowanie bąbelkowe. Drugą najwolniejszą implementacją jest Sortowanie wybieranie,a trzecią Sortowanie wstawianie. Sortowanie scalanie jest drugim najszybszym algorytmem. Sortowanie zliczanie jest najszybszą implementacji. 
	 
	\medskip
	
	Tak zapisany algorytm jest stabilny, ponieważ elementy o tych samych kluczach mają tą samą, jak na początku kolejność. Działa on w czasie proporcjonalnym do liczby elementów w liście, dlatego że każde sortowanie zliczanie działa w czasie złożoności obliczeniowej.
	\medskip
		
	\end{pysub}
	
	
\end{document}